\documentclass[11pt,a4paper]{article}
% Estilo comum dos PDFs individuais da Sprint 1.
\usepackage[utf8]{inputenc}
\usepackage[T1]{fontenc}
\usepackage[brazilian]{babel}
\usepackage[a4paper,margin=2.35cm]{geometry}
\usepackage{xcolor}
\usepackage{hyperref}
\usepackage{enumitem}
\usepackage{listings}
\usepackage{parskip}
\usepackage{needspace}
\definecolor{azul}{RGB}{20,63,112}
\hypersetup{colorlinks=true,linkcolor=azul,urlcolor=azul}
\setlist[itemize]{leftmargin=*,itemsep=3pt}
\lstset{basicstyle=\ttfamily\small,frame=single,breaklines=true,columns=fullflexible}
\sloppy
\newcommand{\codigo}[1]{\texttt{#1}}
\newcommand{\titulo}[2]{\begin{center}{\LARGE\bfseries #1\par}\vspace{4pt}{\large #2\par}\vspace{8pt}\rule{0.8\textwidth}{0.4pt}\par\vspace{5pt}{\small Matheus Moreli Durão Silva -- Matrícula 12621CBS217\par Cibersegurança - Bacharelado - Noturno\par}\end{center}}
\newcommand{\cardfonte}[1]{\par\smallskip\noindent\textit{Referência: }\href{#1}{Trello - enunciado e checklist do card}.}

\begin{document}\titulo{S1\_10 - Listas}{Cadastrar, remover e listar ativos}
\section*{Entregável correto}
O programa deve manter uma lista de ativos e permitir cadastrar, remover e listar. A lista é adequada quando a quantidade de elementos muda enquanto o programa está em execução.
\section*{Como uma lista funciona}
Uma lista usa colchetes: \codigo{ativos = []}. \codigo{append(valor)} acrescenta ao final; \codigo{remove(valor)} remove o primeiro valor igual; \codigo{in} verifica se algo está presente. Índices começam em zero: \codigo{ativos[0]} é o primeiro item; \codigo{ativos[1:3]} é um recorte. Antes de remover, verificar com \codigo{in} evita erro quando o valor não existe.
\begin{lstlisting}
ativos = []
ativos.append("Roteador do laboratorio")
if "Roteador do laboratorio" in ativos:
    ativos.remove("Roteador do laboratorio")
for ativo in ativos:
    print(ativo)
\end{lstlisting}
\section*{Aplicação no projeto}
O \codigo{app.py} guarda vários ativos em uma coleção e usa \codigo{for} para exibir cada um. O projeto final evolui a estrutura: em vez de apenas texto, cada item possui vários dados em dicionários. A ideia de adicionar, procurar e remover continua a mesma.
\section*{Checklist demonstrado}
Listas, criação, índice/fatiamento, adição, remoção, laços e aplicação no inventário. O teste mínimo é cadastrar dois nomes, listar os dois, remover um e confirmar que o outro permanece.
\cardfonte{https://trello.com/c/EurwtLl9}
\end{document}
